\chapter{DF-063 Audit app shell family across real surfaces}\label{record-104}

\noindent{\footnotesize\color{BookMuted}\path{docs/design/DF-063-audit-app-shell-family-across-real-surfaces.md}\par}\smallskip

\begin{RecordMeta}
title: DF-063 Audit app shell family across real surfaces
legend: DF
lane: design
priority: medium
keywords:
  - component-audit
  - dogfood
  - lowerings
  - shell

\end{RecordMeta}

\section{Framing}\label{record-104-framing}

The app shell is the place where page content, global status, command discovery, overlays, and shell-owned chrome meet. \texttt{c\allowbreak{}r\allowbreak{}e\allowbreak{}a\allowbreak{}t\allowbreak{}e\allowbreak{}F\allowbreak{}r\allowbreak{}a\allowbreak{}m\allowbreak{}e\allowbreak{}d\allowbreak{}A\allowbreak{}p\allowbreak{}p\allowbreak{}(\allowbreak{})} owns the high-level shell, while \texttt{s\allowbreak{}t\allowbreak{}a\allowbreak{}t\allowbreak{}u\allowbreak{}s\allowbreak{}B\allowbreak{}a\allowbreak{}r\allowbreak{}(\allowbreak{})} / \texttt{s\allowbreak{}t\allowbreak{}a\allowbreak{}t\allowbreak{}u\allowbreak{}s\allowbreak{}B\allowbreak{}a\allowbreak{}r\allowbreak{}S\allowbreak{}u\allowbreak{}r\allowbreak{}f\allowbreak{}a\allowbreak{}c\allowbreak{}e\allowbreak{}(\allowbreak{})} and \texttt{c\allowbreak{}o\allowbreak{}m\allowbreak{}m\allowbreak{}a\allowbreak{}n\allowbreak{}d\allowbreak{}P\allowbreak{}a\allowbreak{}l\allowbreak{}e\allowbreak{}t\allowbreak{}t\allowbreak{}e\allowbreak{}(\allowbreak{})} / \texttt{c\allowbreak{}o\allowbreak{}m\allowbreak{}m\allowbreak{}a\allowbreak{}n\allowbreak{}d\allowbreak{}P\allowbreak{}a\allowbreak{}l\allowbreak{}e\allowbreak{}t\allowbreak{}t\allowbreak{}e\allowbreak{}S\allowbreak{}u\allowbreak{}r\allowbreak{}f\allowbreak{}a\allowbreak{}c\allowbreak{}e\allowbreak{}(\allowbreak{})} are the smaller seams that make shell state visible without pushing unrelated metadata into page content.

This audit verifies the existing DOGFOOD \texttt{a\allowbreak{}p\allowbreak{}p\allowbreak{}-\allowbreak{}s\allowbreak{}h\allowbreak{}e\allowbreak{}l\allowbreak{}l} story against real story rendering, component-family doctrine, and mode-aware lowerings. It does not redesign the framed runtime; it turns current shell story truth into executable release proof.

\section{Sponsored Users}\label{record-104-sponsored-users}

\begin{itemize}
\tightlist
\item
  TUI app authors building multi-page framed applications with global status and command discovery.
\item
  Docs readers checking how a rich shell lowers when command palettes, status rails, and background interactivity are unavailable.
\item
  Maintainers who need one focused regression tying app-shell story metadata, variants, and component-family guidance to real output.
\end{itemize}

\section{Hills}\label{record-104-hills}

\begin{enumerate}
\def\labelenumi{\arabic{enumi}.}
\tightlist
\item
  A builder can inspect DOGFOOD and see the difference between framed page chrome and command-discovery overlay behavior.
\item
  A reader can switch both app-shell variants through interactive, static, pipe, and accessible profiles while preserving current page, status, and command-discovery facts.
\item
  A maintainer can run one cycle test that proves shell story metadata and design-system guidance have not drifted away from runtime behavior.
\end{enumerate}

\section{Playback Questions}\label{record-104-playback-questions}

\begin{itemize}
\tightlist
\item
  Does DOGFOOD expose the family through the \texttt{a\allowbreak{}p\allowbreak{}p\allowbreak{}-\allowbreak{}s\allowbreak{}h\allowbreak{}e\allowbreak{}l\allowbreak{}l} story?
\item
  Do variants cover both framed-page status chrome and command-discovery overlay behavior?
\item
  Do rich/static previews keep header, footer, page content, and palette geometry visible?
\item
  Do pipe lowerings preserve minimal shell context without pretending hidden regions are still interactive?
\item
  Do accessible lowerings linearize page, status, and command-discovery facts in reading order?
\item
  Do component-family docs still separate shell chrome from page content, notification events, and command execution?
\end{itemize}

\section{Requirements}\label{record-104-requirements}

\begin{itemize}
\tightlist
\item
  Keep the existing \texttt{a\allowbreak{}p\allowbreak{}p\allowbreak{}-\allowbreak{}s\allowbreak{}h\allowbreak{}e\allowbreak{}l\allowbreak{}l} story identity and \texttt{a\allowbreak{}p\allowbreak{}p\allowbreak{}-\allowbreak{}s\allowbreak{}h\allowbreak{}e\allowbreak{}l\allowbreak{}l} coverage family id.
\item
  Render every app-shell story variant in every canonical story profile.
\item
  Treat box drawing and status-rail fill characters as visual-profile chrome only.
\item
  Keep pipe and accessible lowerings shell-context-first and semantic-preserving.
\item
  Confirm \texttt{d\allowbreak{}o\allowbreak{}c\allowbreak{}s\allowbreak{}/\allowbreak{}d\allowbreak{}e\allowbreak{}s\allowbreak{}i\allowbreak{}g\allowbreak{}n\allowbreak{}-\allowbreak{}s\allowbreak{}y\allowbreak{}s\allowbreak{}t\allowbreak{}e\allowbreak{}m\allowbreak{}/\allowbreak{}c\allowbreak{}o\allowbreak{}m\allowbreak{}p\allowbreak{}o\allowbreak{}n\allowbreak{}e\allowbreak{}n\allowbreak{}t\allowbreak{}-\allowbreak{}f\allowbreak{}a\allowbreak{}m\allowbreak{}i\allowbreak{}l\allowbreak{}i\allowbreak{}e\allowbreak{}s\allowbreak{}.\allowbreak{}m\allowbreak{}d} still matches the runtime story posture.
\end{itemize}

\section{Acceptance Criteria}\label{record-104-acceptance-criteria}

\begin{itemize}
\tightlist
\item
  \texttt{t\allowbreak{}e\allowbreak{}s\allowbreak{}t\allowbreak{}s\allowbreak{}/\allowbreak{}c\allowbreak{}y\allowbreak{}c\allowbreak{}l\allowbreak{}e\allowbreak{}s\allowbreak{}/\allowbreak{}D\allowbreak{}F\allowbreak{}-\allowbreak{}0\allowbreak{}6\allowbreak{}3\allowbreak{}/\allowbreak{}a\allowbreak{}p\allowbreak{}p\allowbreak{}-\allowbreak{}s\allowbreak{}h\allowbreak{}e\allowbreak{}l\allowbreak{}l\allowbreak{}-\allowbreak{}f\allowbreak{}a\allowbreak{}m\allowbreak{}i\allowbreak{}l\allowbreak{}y\allowbreak{}-\allowbreak{}a\allowbreak{}u\allowbreak{}d\allowbreak{}i\allowbreak{}t\allowbreak{}.\allowbreak{}t\allowbreak{}e\allowbreak{}s\allowbreak{}t\allowbreak{}.\allowbreak{}t\allowbreak{}s} proves the cycle doc carries the modern playback sections.
\item
  The test locates the DOGFOOD \texttt{c\allowbreak{}r\allowbreak{}e\allowbreak{}a\allowbreak{}t\allowbreak{}e\allowbreak{}F\allowbreak{}r\allowbreak{}a\allowbreak{}m\allowbreak{}e\allowbreak{}d\allowbreak{}A\allowbreak{}p\allowbreak{}p\allowbreak{}(\allowbreak{})\allowbreak{} \allowbreak{}/\allowbreak{} \allowbreak{}s\allowbreak{}t\allowbreak{}a\allowbreak{}t\allowbreak{}u\allowbreak{}s\allowbreak{}B\allowbreak{}a\allowbreak{}r\allowbreak{}S\allowbreak{}u\allowbreak{}r\allowbreak{}f\allowbreak{}a\allowbreak{}c\allowbreak{}e\allowbreak{}(\allowbreak{})\allowbreak{} \allowbreak{}/\allowbreak{} \allowbreak{}c\allowbreak{}o\allowbreak{}m\allowbreak{}m\allowbreak{}a\allowbreak{}n\allowbreak{}d\allowbreak{}P\allowbreak{}a\allowbreak{}l\allowbreak{}e\allowbreak{}t\allowbreak{}t\allowbreak{}e\allowbreak{}S\allowbreak{}u\allowbreak{}r\allowbreak{}f\allowbreak{}a\allowbreak{}c\allowbreak{}e\allowbreak{}(\allowbreak{})} story and verifies expected variants.
\item
  Every variant renders non-empty output in every documented profile.
\item
  Rich/static output includes structured visual containment; pipe and accessible output do not depend on box drawing.
\item
  Pipe and accessible output preserve shell title, current page, status, and command palette intent.
\item
  Component-family docs retain framed app, status bar, command palette, shell-state, pipe, and accessible guidance.
\end{itemize}

\section{Implementation Outline}\label{record-104-implementation-outline}

\begin{enumerate}
\def\labelenumi{\arabic{enumi}.}
\tightlist
\item
  Add a DF-063 cycle test that renders the \texttt{a\allowbreak{}p\allowbreak{}p\allowbreak{}-\allowbreak{}s\allowbreak{}h\allowbreak{}e\allowbreak{}l\allowbreak{}l} DOGFOOD story through the shared story protocol.
\item
  Read-test the app shell section in \texttt{d\allowbreak{}o\allowbreak{}c\allowbreak{}s\allowbreak{}/\allowbreak{}d\allowbreak{}e\allowbreak{}s\allowbreak{}i\allowbreak{}g\allowbreak{}n\allowbreak{}-\allowbreak{}s\allowbreak{}y\allowbreak{}s\allowbreak{}t\allowbreak{}e\allowbreak{}m\allowbreak{}/\allowbreak{}c\allowbreak{}o\allowbreak{}m\allowbreak{}p\allowbreak{}o\allowbreak{}n\allowbreak{}e\allowbreak{}n\allowbreak{}t\allowbreak{}-\allowbreak{}f\allowbreak{}a\allowbreak{}m\allowbreak{}i\allowbreak{}l\allowbreak{}i\allowbreak{}e\allowbreak{}s\allowbreak{}.\allowbreak{}m\allowbreak{}d}.
\item
  Move the backlog note into \texttt{d\allowbreak{}o\allowbreak{}c\allowbreak{}s\allowbreak{}/\allowbreak{}d\allowbreak{}e\allowbreak{}s\allowbreak{}i\allowbreak{}g\allowbreak{}n\allowbreak{}/} and update the v6 lane pointer.
\item
  Record any drift as follow-on backlog debt instead of widening this audit.
\end{enumerate}

\section{Drift Check}\label{record-104-drift-check}

No runtime drift was found in this slice. The existing \texttt{a\allowbreak{}p\allowbreak{}p\allowbreak{}-\allowbreak{}s\allowbreak{}h\allowbreak{}e\allowbreak{}l\allowbreak{}l} story already distinguishes framed page chrome from command discovery, renders structured shell surfaces in interactive/static profiles, and lowers to minimal page/status/palette summaries in pipe/accessibility modes.

No follow-on backlog item is needed for this slice.

\section{Playback}\label{record-104-playback}

\begin{itemize}
\tightlist
\item
  RED: the release lane had only a backlog note; no DF-063 playback test tied app shell behavior to every canonical profile.
\item
  GREEN: the new cycle test renders both \texttt{f\allowbreak{}r\allowbreak{}a\allowbreak{}m\allowbreak{}e\allowbreak{}d\allowbreak{}-\allowbreak{}p\allowbreak{}a\allowbreak{}g\allowbreak{}e} and \texttt{c\allowbreak{}o\allowbreak{}m\allowbreak{}m\allowbreak{}a\allowbreak{}n\allowbreak{}d\allowbreak{}-\allowbreak{}d\allowbreak{}i\allowbreak{}s\allowbreak{}c\allowbreak{}o\allowbreak{}v\allowbreak{}e\allowbreak{}r\allowbreak{}y} variants across interactive, static, pipe, and accessible profiles.
\item
  Rich/static previews preserve header, footer, page, and command-palette geometry.
\item
  Pipe previews preserve current page and shell status without box drawing.
\item
  Accessible previews linearize the same shell facts in reading order.
\end{itemize}

\section{Retrospective}\label{record-104-retrospective}

The app-shell story did not need runtime changes. The useful release work was proving that shell chrome is shared application truth, not incidental drawing inside a page renderer.

